\SourceLeaf{78}\begin{center}
{\Large FOURTH SECTION}\par\bigskip
{\small ON}\par\bigskip
{\Large CONGRUENCES OF THE SECOND DEGREE.}
\end{center}

\begin{center}\rule{.09\linewidth}{.25pt}\end{center}

\Needspace{10\baselineskip}
\begin{center}\small\emph{Quadratic residues and nonresidues.}\end{center}

\Article{94}

{\rightskip=0pt plus 3em
\textsc{Theorem.} \emph{For any number $m$ taken as modulus, among the numbers \mbox{$0,1,2,3....m-1$}, at most $\tfrac12m+1$ can be congruent to a square when $m$ is even, and at most $\tfrac12m+\tfrac12$ when $m$ is odd.}\par}

\emph{Proof.} Since the squares of congruent numbers are congruent, any number that can be congruent to a square will also be congruent to a square whose root is $<m$. So it is enough to consider the least residues of the squares $0,1,4,9....(m-1)^2$. But we readily see that $(m-1)^2\equiv1$, $(m-2)^2\equiv2^2$, $(m-3)^2\equiv3^2$, etc. Thus, when $m$ is even, the squares $(\tfrac12m-1)^2$ and $(\tfrac12m+1)^2$, $(\tfrac12m-2)^2$ and $(\tfrac12m+2)^2$, etc. have the same least residues; when $m$ is odd, the squares $(\tfrac12m-\tfrac12)^2$ and $(\tfrac12m+\tfrac12)^2$, $(\tfrac12m-\tfrac32)^2$ and $(\tfrac12m+\tfrac32)^2$, etc. are congruent. So no numbers other than those congruent to one of the squares $0,1,4,9....(\tfrac12m)^2$ can be congruent to a square when $m$ is even; when it is odd, every number congruent to a square must be congruent to one of $0,1,4,9....(\tfrac12m-\tfrac12)^2$. There are therefore at most $\tfrac12m+1$ distinct least residues in the first case, and $\tfrac12m+\tfrac12$ in the second. \mbox{\emph{Q. E. D.}}

\emph{Example.} Modulo $13$, the squares of the numbers $0,1,2,3...6$ have the least residues $0,1,4,9,3,12,10$; after these, the same residues \SourceLeaf{79}recur in reverse order, $10,12,3$, etc. So any number congruent to none of these residues, or equivalently congruent to one of $2,5,6,7.\ 8,11$, cannot be congruent to a square.

Modulo $15$ we find the residues $0,1,4,9,1,10,6,4$, after which the same residues recur in reverse order. Here the number of residues that can be congruent to a square is even smaller than $\tfrac12m+\tfrac12$, since they are $0,1,4,6,9,10$. The numbers $2,3,5,7,8,11,12,13,14$, and those congruent to any one of them, cannot be congruent to a square modulo $15$.

\Article{95}
It follows that, for any modulus, all numbers can be divided into two classes: one contains the numbers that can be congruent to a square, and the other those that cannot. We will call the former \emph{quadratic residues of the number taken as modulus}\HistoricalNote{*)}{Strictly speaking, we are using the genitive here in a different sense from before. We should say that $r$ is a residue of the square $aa$ modulo $m$ when $r\equiv aa\,(\operatorname{mod.}m)$; but for brevity, throughout this section we call \emph{$r$} a quadratic residue \emph{of $m$ itself}, without any risk of ambiguity. From now on we will not use the word \emph{residue} to mean a congruent number, except when speaking of \emph{least} residues, where no doubt can arise.}, and the latter its \emph{quadratic nonresidues}, or simply \emph{residues} and \emph{nonresidues} whenever there is no risk of ambiguity. It is enough to classify all the numbers $0,1,2\ldots m-1$, since congruent numbers must belong to the same class.

In this investigation too, we will begin with prime moduli; this is to be understood even when we do not say so explicitly. The prime number $2$ must, however, be excluded: we will consider only \emph{odd} primes.

\Needspace{10\baselineskip}
\begin{center}\small\emph{When the modulus is prime, there are as many residues smaller than it as nonresidues.}\end{center}

\Article{96}
\emph{Taking the prime number $p$ as modulus, half the numbers $1,2,3\GaussFourDots p-1$ are quadratic residues, and the rest are nonresidues; that is, there are $\tfrac12(p-1)$ residues and just as many nonresidues.}

For it is easy to prove that all the squares $1,4,9\ldots\tfrac14(p-1)^2$ are incongruent. If $rr\equiv r'r'(\operatorname{mod.}p)$ were possible with the numbers $r,r'$ unequal and no greater than $\tfrac12(p-1)$, then, taking $r>r'$, as we may, $(r-r')(r+r')$ would be positive and \SourceLeaf{80}divisible by $p$. But both factors $r-r'$ and $r+r'$ are smaller than $p$, so this assumption is impossible (Article 13). Thus there are $\tfrac12(p-1)$ quadratic residues among the numbers $1,2,3\GaussFourDots p-1$. There cannot be more, since adding the residue $0$ gives $\tfrac12(p+1)$, and the total number of residues cannot exceed this. So the remaining numbers are nonresidues, and their number is $=\tfrac12(p-1)$.

Since zero is always a residue, we exclude it and numbers divisible by the modulus from these investigations. This case needs no explanation and would only spoil the neatness of the theorems. For the same reason we have excluded the modulus $2$.

\Article{97}
Several results that we will present in this section can also be derived from the principles of the previous section. Since it is useful to examine the same truth by different methods, we will show this connection. We readily see that all numbers congruent to a square have \emph{even} indices, while those that cannot be congruent to any square have odd indices. Since $p-1$ is even, there are as many even indices as odd ones, namely $\tfrac12(p-1)$, and therefore as many residues as nonresidues.

\emph{Examples.}\quad The residues for the following moduli are
\begin{center}
\begin{tabular}{@{}r@{\quad}c@{\quad}l@{}}
$3$ & $\mathinner{\ldotp\ldotp\ldotp\ldotp\ldotp}$ & $1.$\\
$5$ & $\mathinner{\ldotp\ldotp\ldotp\ldotp\ldotp}$ & $1,4.$\\
$7$ & $\mathinner{\ldotp\ldotp\ldotp\ldotp\ldotp}$ & $1,2,4.$\\
$11$ & $\mathinner{\ldotp\ldotp\ldotp\ldotp\ldotp}$ & $1,3,4,5,9.$\\
$13$ & $\mathinner{\ldotp\ldotp\ldotp\ldotp\ldotp}$ & $1.3.4,9,10,12.$\\
$17$ & $\GaussFourDots,$ & $1,2,4,8,9,13,15,16.$\\
\multicolumn{3}{c}{etc.}
\end{tabular}
\end{center}
The other numbers smaller than these moduli are nonresidues.

\Needspace{10\baselineskip}
\begin{center}\small\emph{Whether a composite number is a residue or a nonresidue of a given prime depends on its factors.}\end{center}

\Article{98}
\textsc{Theorem.}\quad \emph{The product of two quadratic residues of the prime number $p$ is a residue; the product of a residue and a nonresidue is a nonresidue; and the product of two nonresidues is a residue.}

\SourceLeaf{81}\emph{Proof.}\quad I.\quad Let $A,B$ be residues arising from the squares $aa,bb$, so that $A\equiv aa$, $B\equiv bb$. Their product $AB$ is congruent to the square of $ab$, and is therefore a residue.

II.\quad When $A$ is a residue, say $\equiv aa$, and $B$ is a nonresidue, $AB$ is a nonresidue. For suppose, if possible, that $AB\equiv kk$, and let $\frac{k}{a}\,(\mathrm{mod.}\,p)\equiv b$. Then $aaB\equiv aabb$, so $B\equiv bb$: that is, $B$ is a residue, contrary to the hypothesis.

\emph{Another proof.}\quad Multiply all the residues among $1,2,3\GaussFourDots p-1$ (whose number is $=\tfrac12(p-1)$) by $A$. All the products are quadratic residues, and they are all incongruent. If we now multiply the nonresidue $B$ by $A$, the product will be congruent to none of those already obtained. If it were a residue, there would therefore be $\tfrac12(p+1)$ incongruent residues, not yet including the residue $0$, contrary to Article 96.

III.\quad Let $A,B$ be nonresidues. Multiply all the residues among $1,2,3\GaussFourDots p-1$ by $A$; we obtain $\tfrac12(p-1)$ mutually incongruent nonresidues (II). The product $AB$ cannot be congruent to any of them; so if it were a nonresidue, there would be $\tfrac12(p+1)$ mutually incongruent nonresidues, contrary to Article 96. So the product, etc. \mbox{\emph{Q. E. D.}}

These theorems follow still more easily from the principles of the previous section. The indices of residues are always even, and those of nonresidues odd; so the index of the product of two residues or of two nonresidues is even, and the product itself is a residue. In contrast, the index of the product of a residue and a nonresidue is odd, so the product itself is a nonresidue.

Both methods of proof can also be used for the following theorems: \emph{The value of the expression $\frac{a}{b}\,(\mathit{mod.}\,p)$ is a residue when the numbers $a,b$ are both residues or both nonresidues; it is a nonresidue when one of the numbers $a,b$ is a residue and the other a nonresidue.} These also follow by reversing the preceding theorems.

\Article{99}
In general, a product of any number of factors is a residue both when all the factors are residues and when an even number of them are nonresidues. When an odd number of the factors are nonresidues, the product is a nonresidue. So we can readily decide \SourceLeaf{82}whether a composite number is a residue if we know the character of each of its factors. For this reason, we have included only prime numbers in Table II. The table is arranged as follows. The moduli are placed in the margin\HistoricalNote{*)}{We will soon show how we can also dispense with composite moduli.}, and the successive prime numbers in the body of the table. When one of these primes is a residue of a modulus, a short line is placed in the corresponding space; when the prime is a nonresidue of the modulus, that space is left empty.

\Needspace{10\baselineskip}
\begin{center}\small\emph{Composite moduli.}\end{center}

\Article{100}
Before moving on to more difficult matters, we must add a few observations about nonprime moduli.

If, for a prime number $p$, a power $p^n$ is taken as modulus (where we assume that $p$ is not $2$), half the numbers not divisible by $p$ and smaller than the modulus are residues, and half are nonresidues; that is, the number of each is $=\tfrac12(p-1)p^{n-1}$.

For if $r$ is a residue, it is congruent to a square whose root does not exceed half the modulus; see Article 94. We readily see that there are $\tfrac12(p-1)p^{n-1}$ numbers not divisible by $p$ and smaller than half the modulus. It remains to prove that their squares are all incongruent, or give distinct quadratic residues. If the squares of two numbers $a,b$, not divisible by $p$ and smaller than half the modulus, were congruent, then $aa-bb$, or $(a-b)(a+b)$, would be divisible by $p^n$ (taking $a>b$, as we may). This is possible only if \emph{either} one of the numbers $a-b.\ a+b$ is divisible by $p^n$, which is impossible since each is $<p^n$, \emph{or} one is divisible by $p^m$ and the other by $p^{n-m}$, so both are divisible by $p$. But this too is impossible. Their sum and difference, $2a$ and $2b$, would then also be divisible by $p$, and therefore so would $a$ and $b$, contrary to the hypothesis. --- We conclude that, among the numbers not divisible by $p$ and smaller than the modulus, there are $\tfrac12(p-1)p^{n-1}$ residues; the remaining numbers, of which there are just as many, are nonresidues. \mbox{\emph{Q. E. D.}} --- This theorem can also be derived by considering indices, much as in Article 97.

\Needspace{8\baselineskip}
\Article{101}
\emph{Every number not divisible by $p$ that is a residue of $p$ is also a residue \SourceLeaf{83}of $p^n$; every nonresidue of $p$ is also a nonresidue of $p^n$.}

The latter part of this proposition needs no proof. If the first part were false, then among the numbers smaller than $p^n$ and not divisible by $p$, more would be residues of $p$ than of $p^n$, that is, more than $\frac{1}{2}p^{n-1}(p-1)$. But we readily see that the number of residues of $p$ among these numbers is exactly $=\frac{1}{2}p^{n-1}(p-1)$.

It is just as easy to find an actual square congruent to a given residue modulo $p^n$, if a square congruent to it modulo $p$ is known.

If we have a square $aa$ congruent to the given residue $A$ modulo $p^\mu$, we can derive a square congruent to $A$ modulo $p^\nu$ (where we assume $\nu>\mu$ and $=$ or $<2\mu$) as follows. Put the root of the required square $=\pm a+xp^\mu$, a form that it must clearly have. We must then have $aa\pm2axp^\mu+xxp^{2\mu}\equiv A\,(\mathrm{mod.}\,p^\nu)$, or, since $2\mu\geq\nu$, $A-aa\equiv\pm2axp^\mu\,(\mathrm{mod.}\,p^\nu)$. Let $A-aa=p^{\mu}d$; then $x$ is the value of the expression $\pm\frac{d}{2a}\,(\mathrm{mod.}\,p^{\nu-\mu})$, which is equivalent to $\pm\frac{A-aa}{2ap^\mu}\,(\mathrm{mod.}\,p^\nu)$.

So a square congruent to $A$ modulo $p$ gives a square congruent to $A$ modulo $p^2$; from here we can move up to the modulus $p^4$, then to $p^8$, etc.

\emph{Example.} Given the residue $6$, congruent modulo $5$ to the square $1$, we find the square $9^2$ to which it is congruent modulo $25$, then $16^2$ to which it is congruent modulo $125$, etc.

\Needspace{7\baselineskip}
\Article{102}
As for numbers divisible by $p$, their squares are clearly divisible by $pp$. So all numbers divisible by $p$ but not by $pp$ are nonresidues of $p^n$. In general, given a number $p^k A$ where $A$ is not divisible by $p$, we must distinguish the following cases:

1) When $k=$ or $>n$, we have $p^k A\equiv0\,(\mathrm{mod.}\,p^n)$, a residue.

2) When $k<n$ and the exponent is odd, $p^k A$ is a nonresidue.

For if $p^k A=p^{2\kappa+1}A\equiv ss\,(\mathrm{mod.}\,p^n)$, then $ss$ would be divisible by $p^{2\kappa+1}$, which is possible only if $s$ is divisible by $p^{\kappa+1}$. But then $ss$ \SourceLeaf{84}would also be divisible by $p^{2\kappa+2}$, and therefore (since $2\kappa+2$ is certainly no greater than $n$) so would $p^k A$, \emph{that is,} $p^{2\kappa+1}A$. Thus $A$ would be divisible by $p$, contrary to the hypothesis.

3) When $k<n$ and the exponent is even. Then $p^k A$ is a residue or a nonresidue of $p^n$ according as $A$ is a residue or a nonresidue of $p$. For when $A$ is a residue of $p$, it is also a residue of $p^{n-k}$. If $A\equiv aa (\mathrm{mod.}\,p^{n-k})$, then $Ap^k\equiv aap^k\,(\mathrm{mod.}\,p^n)$, and $aap^k$ is a square. But when $A$ is a nonresidue of $p$, $p^k A$ cannot be a residue of $p^n$. For suppose $p^k A\equiv aa\,(\mathrm{mod.}\,p^n)$; then $aa$ must be divisible by $p^k$. The quotient is a square congruent to $A$ modulo $p^{n-k}$, and therefore also modulo $p$; \emph{that is,} $A$ would be a residue of $p$, contrary to the hypothesis.

\Article{103}
Since we excluded the case $p=2$, we must now say something about it. When the modulus is $2$, every number is a residue, and there are no nonresidues. When the modulus is $4$, all odd numbers of the form $4k+1$ are residues, and all those of the form $4k+3$ are nonresidues. Finally, when the modulus is $8$ or a higher power of $2$, all odd numbers of the form $8k+1$ are residues, and the others, those of the forms $8k+3,\ 8k+5,\ 8k+7$, are nonresidues. The latter part of this proposition is clear, since the square of any odd number, whether of the form $4k+1$ or $4k-1$, has the form $8k+1$. We prove the first part as follows.

1) If either the sum or the difference of two numbers is divisible by $2^{n-1}$, their squares are congruent modulo $2^n$. For if one is put $=a$, the other has the form $2^{n-1}h\pm a$, whose square is $\equiv aa\,(\mathrm{mod.}\,2^n)$.

2) Every odd number that is a quadratic residue of $2^n$ is congruent to a square whose root is odd and $<2^{n-2}$. Let $aa$ be any square congruent to that number, and let $a\equiv\pm\alpha (\mathrm{mod.}\,2^{n-1})$, with $\alpha$ no greater than half the modulus (Article 4). Then $aa\equiv\alpha\alpha$, so the given number is also $\equiv\alpha\alpha$. Both $a$ and $\alpha$ are clearly odd, and $\alpha<2^{n-2}$.

3) The squares of all odd numbers smaller than $2^{n-2}$ are incongruent modulo $2^n$. For let $r$ and $s$ be two such numbers. If their squares were congruent modulo $2^n$, then $(r-s)(r+s)$ would be divisible by $2^n$ (taking $r>s$). But the numbers $r-s,\ r+s$ clearly cannot both be divisible by $4$; \SourceLeaf{85}so if one is divisible only by $2$, to make their product divisible by $2^n$, the other would have to be divisible by $2^{n-1}$. \mbox{\emph{Q. E. A.}} since both are $<2^{n-2}$.

\Needspace{8\baselineskip}
4) Finally, reduce these squares to their \emph{least positive residues}. We obtain $2^{n-3}$ distinct quadratic residues smaller than the modulus\HistoricalNote{*)}{Since the number of odd numbers below $2^{n-2}$ is $2^{n-3}$.}, each of the form $8k+1$. But there are exactly $2^{n-3}$ numbers of the form $8k+1$ smaller than the modulus, so they must all occur among those residues. \mbox{\emph{Q. E. D.}}

\Needspace{3\baselineskip}
To find a square congruent to a given number of the form $8k+1$ modulo $2^n$, we can use a method like that of Article 101; see also Article 88. --- Finally, what we said in general in Article 102 also applies to even numbers.

\Article{104}
The following facts about the number of distinct values (\emph{that is,} values incongruent modulo the given modulus) admitted by an expression such as $V=\surd A\,(\mathrm{mod.}\,p^n)$, when $A$ is a residue of $p^n$, follow easily from what precedes. (As before, we assume that $p$ is prime, and for brevity include the case $n=1$ from the outset.) I. If $A$ is not divisible by $p$, $V$ has \emph{one} value for $p=2,\ n=1$, namely $V \equiv1$; \emph{two} when $p$ is odd, and also for $p=2,\ n=2$: if one is put $\equiv v$, the other is $\equiv-v$; \emph{four} for $p=2,\ n>2$: if one is put $\equiv v$, the others are $\equiv-v,\ 2^{n-1}+v,\ 2^{n-1}-v$. II. If $A$ is divisible by $p$ but not by $p^n$, let the highest power of $p$ dividing $A$ be $p^{2\mu}$ (its exponent must clearly be even), and let $A=ap^{2\mu}$. All values of $V$ are then clearly divisible by $p^\mu$, and the quotients are values of the expression $V'=\surd a\,(\mathrm{mod.}\,p^{n-2\mu})$. So all distinct values of $V$ are obtained by multiplying all values of the expression $V'$ lying between $0$ and $p^{n-\mu}$ by $p^\mu$. They are therefore given by

\[
vp^\mu,\quad
vp^\mu+p^{n-\mu},\quad
vp^\mu+2p^{n-\mu}\ldots vp^\mu+(p^\mu-1)p^{n-\mu}
\]

where $v$ stands in turn for all \emph{distinct} values of the expression $V'$. Their number is thus $p^\mu,\ 2p^\mu$ or $4p^\mu$, according as the number of the latter values (by case I) is $1,\ 2$ or $4$. III. If $A$ is divisible by $p^n$, put $n=2m$ or $=2m-1$ according as it is even or odd. We readily see that all numbers divisible by $p^m$, and no others, are values of $V$. So all its distinct values are $0,p^m,2p^m\ldots(p^{n-m}-1)p^m$, and their number is $p^{n-m}$.

\SourceLeaf{86}\Article{105}
There remains the case where the modulus $m$ is composed of several primes. Let $m=abc\ldots$, where $a,b,c$, etc. denote distinct primes or powers of distinct primes. If $n$ is a residue of $m$, then $n$ is clearly also a residue of each of $a,b,c$, etc.; so $n$ is certainly a nonresidue of $m$ if it is a nonresidue of any of $a,b,c$, etc. Conversely, if $n$ is a residue of each of $a,b,c$, etc., it is also a residue of their product $m$. For suppose $n\equiv A^2$, $B^2,C^2$, etc. modulo $a,b,c$, etc., respectively. Find a number $N$ congruent to $A,B,C$, etc. modulo $a,b,c$, etc., respectively (Article 32). Then $n\equiv NN$ modulo all these moduli, and therefore modulo their product $m$. --- We readily see that combining \emph{any} value of $A$, or of the expression $\surd n\,(\mathrm{mod.}\,a)$, with \emph{any} value of $B$, with \emph{any} value of $C$, etc., gives a value of $N$, or of the expression $\surd n\,(\mathrm{mod.}\,m)$; that distinct combinations give distinct values of $N$; and that all the combinations give all the values. So the number of all distinct values of $N$ equals the product of the numbers of values of $A,B,C$, etc., which we showed how to determine in the preceding article. --- Furthermore, if one value of the expression $\surd n\,(\mathrm{mod.}\,m)$, or of $N$, is known, it is also a value of all of $A,B,C$, etc. Since all the other values of these quantities can then be derived by the preceding article, all the other values of $N$ can readily be obtained from one.

\Needspace{6\baselineskip}
\emph{Example.} Let the modulus be $315$, and ask whether $46$ is a residue or a nonresidue of it. The prime divisors of $315$ are $3,5,7$, and $46$ is a residue of each of them, so it is also a residue of $315$. Furthermore, since $46\equiv1$ and $\equiv64 (\mathrm{mod.}\,9)$; $\equiv1$ and $\equiv16\,(\mathrm{mod.}\,5)$; $\equiv4$ and $\equiv25\,(\mathrm{mod.}\,7)$, we find the roots of the squares to which $46$ is congruent modulo $315$: $19,26,44,89,226$, $271,289,296$.

\Needspace{10\baselineskip}
\begin{center}
\small\itshape
A general criterion for deciding whether a given number is a residue or a nonresidue of a given prime.
\end{center}

\Article{106}
It follows from what precedes that, if we can always decide whether a given \emph{prime number} is a residue or a nonresidue of a \emph{given prime}, all other cases can be reduced to this one. We must therefore make every effort to find reliable criteria for this case. Before starting that investigation, we give a criterion drawn from the previous section which, though \SourceLeaf{87}of almost no practical use, deserves mention for its simplicity and generality.

\emph{Any number $A$ not divisible by the prime $2m+1$ is a residue or a nonresidue of that prime according as $A^m\equiv+1$ or $\equiv-1\,(\operatorname{mod.}2m+1)$.}

For the modulus $2m+1$, let the index of $A$ in any system be $a$. Then $a$ is even when $A$ is a residue of $2m+1$, and odd when $A$ is a nonresidue. The index of $A^m$ is $ma$, \emph{that is,} $\equiv0$ or $\equiv m\,(\operatorname{mod.}2m)$, according as $a$ is even or odd. Thus $A^m$ is $\equiv+1$ in the first case, and $\equiv-1\,(\operatorname{mod.}2m+1)$ in the second. See Articles 57 and 62.

\emph{Example.}\quad $3$ is a residue of $13$ because $3^6\equiv1\,(\operatorname{mod.}13)$, while $2$ is a nonresidue of $13$ because $2^6\equiv-1\,(\operatorname{mod.}13)$.

But whenever the numbers to be tested are even moderately large, this criterion is quite useless because of the enormous amount of calculation.

\Needspace{10\baselineskip}
\begin{center}\emph{Investigations of the primes for which given numbers are residues or nonresidues.}\end{center}

\Article{107}
Given a modulus, it is very easy to specify all the numbers that are its residues or nonresidues. Put that number $=m$. We must determine the squares whose roots do not exceed half of $m$, or numbers congruent to these squares modulo $m$ (still quicker methods are available in practice). Then all numbers congruent to one of them modulo $m$ are residues of $m$, and all numbers congruent to none of them are nonresidues. --- But the inverse question, \emph{given a number, to specify all numbers of which it is a residue or a nonresidue}, calls for a much deeper investigation. We will therefore examine this problem, on whose solution the one proposed in the preceding article depends, starting with the simplest cases.

\Needspace{10\baselineskip}
\begin{center}\emph{The residue $-1$.}\end{center}

\Article{108}
\textsc{Theorem.}\quad \emph{For every prime of the form $4n+1$, $-1$ is a quadratic residue; for every prime of the form $4n+3$, it is a nonresidue.}

\emph{Example.}\quad $-1$ is a residue of the numbers $5,\allowbreak 13,\allowbreak 17,\allowbreak 29,\allowbreak 37,\allowbreak 41,\allowbreak 53,\allowbreak 61.73,\allowbreak 89,\allowbreak 97$, etc., arising respectively \SourceLeaf{88}from the squares of $2,\allowbreak 5,\allowbreak 4,\allowbreak 12,\allowbreak 6,\allowbreak 9,\allowbreak 23,\allowbreak 11,\allowbreak 27,\allowbreak 34,\allowbreak 22$, etc.; it is a nonresidue of the numbers $3,\allowbreak 7,\allowbreak 11,\allowbreak 19,\allowbreak 23,\allowbreak 31,\allowbreak 43,\allowbreak 47,\allowbreak 59,\allowbreak 67,\allowbreak 71,\allowbreak 79,\allowbreak 83$, etc.

We already mentioned this theorem in Article 64. Its proof follows easily from Article 106: for a prime of the form $4n+1$ we have $(-1)^{2n}\equiv1$, and for one of the form $4n+3$ we have $(-1)^{2n+1}\equiv-1$. This proof agrees with the one given there. But the theorem is so elegant and useful that another proof will not be out of place.

\Article{109}
Denote the collection of all residues of the prime $p$ that are smaller than $p$, excluding the residue $0$, by the letter $C$. Their number is always $=\frac{p-1}{2}$, so it is even when $p$ has the form $4n+1$, and odd when $p$ has the form $4n+3$. As in Article 77, where we considered numbers in general, call residues \emph{partners} when their product is $\equiv1$ (mod.~$p$); for if $r$ is a residue, $\frac{1}{r}$ (mod.~$p$) is clearly a residue too. Since a residue cannot have more than one partner among the residues $C$, all the residues $C$ can be divided into classes, each containing two partner residues. If no residue were its own partner, \emph{that is,} if each class contained two \emph{unequal} residues, the total number of residues would be twice the number of classes. But if some residues are their own partners, \emph{that is,} if some classes contain only one residue, or, if preferred, the same residue twice, let the number of these classes be $=a$, and that of the remaining classes $=b$; then the number of all residues $C$ is $=a+2b$. So when $p$ has the form $4n+1$, $a$ is even; when $p$ has the form $4n+3$, $a$ is odd. But no numbers smaller than $p$ other than $1$ and $p-1$ can be their own partners (see Article 77), and the first, $1$, certainly occurs among the residues. So in the first case $p-1$ (or, equivalently here, $-1$) must be a residue, and in the second a nonresidue; otherwise in the former case we would have $a=1$, and in the latter $=2$, which is impossible.

\Article{110}
This proof too is due to the distinguished Euler, who also first found the earlier one. See \emph{Opusc. Anal. Vol. I, p. 135.} --- It is easy to see that it rests on principles like those of our second proof of Wilson's theorem in Article 77. But \SourceLeaf{89}if we assume that theorem, we can give an even easier proof. Among the numbers $1,2,3\ldots p-1$ there are $\frac{p-1}{2}$ quadratic residues of $p$ and just as many nonresidues; so the number of nonresidues is even when $p$ has the form $4n+1$, and odd when $p$ has the form $4n+3$. Thus the product of all the numbers $1,2,3\ldots p-1$ is a residue in the first case and a nonresidue in the second (Article 99). But this product is always $\equiv-1$ (mod.~$p$); so $-1$ too is a residue in the first case and a nonresidue in the second.

\Needspace{12\baselineskip}
\Article{111}
So if $r$ is a residue of a prime of the form $4n+1$, $-r$ is also a residue of that prime, while all the nonresidues of such a prime remain nonresidues when their signs are reversed\HistoricalNote{*)}{So when we speak of any number as a residue or a nonresidue of a number of the form $4n+1$, we can disregard its sign altogether, or even give it the double sign $\pm$.}. The opposite happens for primes of the form $4n+3$: reversing the sign turns residues into nonresidues and vice versa. See Article 98.

We also readily derive the general rule: $-1$ \emph{is a residue of every number divisible neither by $4$ nor by any prime of the form $4n+3$, and a nonresidue of every other number.} See Articles 103 and 105.
